\section{Introduction}
\label{sec:intro}

Neuro-symbolic (NeSy) AI sets out to build a single system that \emph{perceives} with neural networks
and \emph{reasons} with symbolic structure~\citep{garcez2020,marra2021}. The appeal is concrete: by
factoring a task into ``recognize the parts'' and ``combine them with known rules,'' a NeSy model
should learn from far less data, expose interpretable intermediate symbols, and reuse those symbols
to solve problems it was never trained on. The canonical demonstration is \mnistadd
\citep{manhaeve2018}: a model sees two digit images and is supervised only on their \emph{sum}, never
on the digits, yet is expected to learn to read digits as a by-product of learning to add.

\para{The promise hides an assumption.} Every downstream NeSy benefit---interpretability, knowledge
reuse, compositional generalization---depends on the latent symbols being \emph{correct}. If the
network learns a mapping from images to symbols that is systematically wrong but still produces the
right final answer, the system has not understood anything; it has found a \newterm{reasoning
shortcut}~\citep{li2024softened,deepsymbolic2022}. A practitioner who watches only task accuracy
will never notice. This is not a hypothetical: with only distant supervision, the perception-to-symbol
map is under-constrained, and training can collapse onto a spurious but self-consistent solution.

\para{What is missing.} The two headline NeSy claims---data efficiency and correct
grounding---are widely asserted but rarely measured together with statistical rigor on identical
splits. Data efficiency is seldom reported head-to-head against a neural baseline with confidence
intervals, and \emph{grounding accuracy is almost never reported at all}, so reasoning shortcuts stay
invisible. Prior work raises the shortcut problem qualitatively, but a controlled, multi-seed
\emph{measurement} of \emph{when} shortcuts appear, \emph{what they look like}, and \emph{which cheap
fixes work} has been missing.

\para{Our approach.} We run a single controlled study on \mnistadd that measures four quantities
side by side, on identical splits, across $5$ seeds, with Welch $t$-tests, Holm--Bonferroni
correction, and Cohen's $d$ effect sizes. We compare three models on a \emph{shared} CNN encoder so
the only variable is the learning signal: a \neural baseline that predicts the sum directly, a
\nesy model that outputs per-image digit distributions and computes the sum by exact differentiable
marginalization (the DeepProbLog/Scallop idea in compact form), and an \oracle trained on true digit
labels as an upper bound. The study is held together by one clean manipulation: \emph{flipping a single
symbolic operator}, standard addition $a{+}b$ versus modular addition \modten, which turns reasoning
shortcuts on and off while everything else---architecture, data, optimizer---stays fixed.

\para{What we find.} (i) Integration is a large, statistically significant data-efficiency win:
at $1{,}000$ pairs the \nesy model reaches $85.4\%$ sum accuracy versus the neural baseline's
$22.0\%$ (Cohen's $d{=}56$, $p{=}2.9{\times}10^{-13}$), recovering digits at $92$--$99\%$ with no
digit labels. (ii) Yet under the modular operator, correct grounding \emph{never} emerges: among
seeds that learn the task, raw grounding is $\approx0.1\%$, while a permutation-aligned metric reveals
a clean systematic relabeling---sometimes the exact cyclic shift $d\mapsto(d{+}5)\bmod 10$. The
oracle grounds at $99\%$ on the \emph{same} task, so the failure is caused by distant supervision under
a symmetric operator, not by the task being hard. (iii) Just $10$ labeled images break the symmetry
and restore grounding to $97.7\%$ in every seed; unsupervised regularizers do not. (iv) With correct
grounding, the \nesy model solves $2$- and $3$-digit addition it never saw at $95.7\%$/$93.7\%$ by
composing its digit classifier with a symbolic place-value sum, while the neural baseline scores
exactly $0\%$.

\para{Contributions.}
\begin{itemize}[leftmargin=*,itemsep=0pt,topsep=0pt]
    \item We quantify NeSy data efficiency head-to-head against a neural baseline and a digit-supervised
    oracle on identical splits, with confidence intervals and effect sizes at six training sizes.
    \item We introduce a permutation-aware (Hungarian-aligned) grounding diagnostic and use a
    single-operator flip ($a{+}b$ vs.\ \modten) to show that whether high task accuracy implies correct
    symbols is a property of the \emph{task's symmetry}, not of NeSy in general.
    \item We compare cheap mitigations and show that a handful of labeled symbols ($10$ images) is the
    decisive, reliable fix for a detected shortcut, while unsupervised priors are not.
    \item We demonstrate that correct grounding unlocks zero-shot compositional generalization to
    multi-digit addition---a capability the neural baseline structurally lacks.
\end{itemize}

The rest of the paper reviews related work (\secref{sec:related}), details our models, operators, and
metrics (\secref{sec:method}), presents results E1--E5 (\secref{sec:results}), discusses
interpretation and limitations (\secref{sec:discussion}), and concludes (\secref{sec:conclusion}).
